\section{Theory}
\label{app:theory}

In this section, we state the results discussed in \cref{sec:signflips}, together with the proof of \cref{thm:unexplained_ordering_sensitivity} and an example showing that the difference between the two decompositions can be made arbitrarily large even when the mean gap is fixed.

%%% MOVED here from \cref{sec:signflips} (main text).

\subsection{Characterizing sign flips}
\label{sec:sign_flip_characterizations}

%%% AS: the sentences below felt redundant with the intro to this section
% We start by concretizing the type of decision we will focus on. 
% In particular, we will consider a conclusion based on the sign of either the explained or unexplained component of the OBD.
% We will assume we are working with full population information so that a sign flip alone is enough to change conclusions.

% Typically the OBD represents the special case of \cref{eq:OB_H_Counterfactual,eq:OB_K_Counterfactual} 
% where the conditional expectation of the outcome given the covariate is linear, as we describe next.
% We maintain focus on linear models for the following reasons:
% (1) The linear model version of the OBD remains the one widely used in practice in the present day.
% (2) Extensions involving penalized regression (e.g., lasso) maintain the same linear-in-features structure.
% (3) Linearity gives us a tractable way of measuring model complexity in terms of the number of included features.
% As we identify substantive problems with correctly specified linear models; we expect practical issues only compound beyond this relatively simple case.

%Our goal is to study conditions on population-level quantities that produce sign flips in the explained and unexplained components.
%For tractability, we begin by assuming finite-dimensional covariates, $X \in \mathbb{R}^d$, and a linear model for $\mathbb{E}_g[Y \mid X]$: $\alpha_g + X^\top \beta_g$
%for some intercept $\alpha_g \in \mathbb{R}$ and slope $\beta_g \in \mathbb{R}^d$.
%\textcolor{red}{Throughout, we will assume an infinite data limit}, where $\alpha_g, \beta_g$ are the parameters of the population-level best linear fit to $\mathbb{E}_g[Y\mid X]$.
%\textcolor{red}{In addition to being theoretically tractable, the linearity assumption is practically relevant, as the OBD is often estimated using OLS to model $\mathbb{E}_g[Y\mid X]$, and OLS coefficients converge to $\alpha_g, \beta_g$ if an appropriate law of large numbers applies.}
%Then, for groups $g,g'$, we observe $\E_{g'} [ \E_g  [ Y | X ] ] = \alpha_g + \E_{g'}[X]^\top \beta_g$.

%We introduce the notation
%$\mu_g := \mathbb{E}_g[X],
%\Delta \mu := \mu_H - \mu_K,
%\Delta \beta := \beta_H - \beta_K,
%\Delta \alpha := \alpha_H - \alpha_K.
%$
%With these choices, the OBD
%(\cref{eq:OB_H_Counterfactual,eq:OB_K_Counterfactual}) becomes
%\begin{equation}
%    \Delta_Y = \underbrace{\Delta \mu^\top \beta_H}_{E^{(H)}} + \underbrace{\mu_K^\top \Delta \beta + \Delta \alpha}_{U^{(H)}} = \underbrace{\Delta \mu^\top \beta_K}_{E^{(K)}} + \underbrace{\mu_H^\top \Delta \beta + \Delta \alpha}_{U^{(K)}}.
%\end{equation}
%The first equality is the OBD with reference group $H$, and the second with reference group $K$.

We start by characterizing sign flips in the explained and unexplained components, respectively; these characterizations will help us later establish
that sign flips of either component are common in parameter space.

\textbf{Explained component.} Explained component sign flips are relatively straightforward to characterize.
%We now study when each of the explained component can flip sign.
\begin{definition}
	\label{defn:sign_flip_explained}
	We say the OBD \emph{explained component flips signs} if the sign of the explained term, $\Delta\mu^T \beta_g$, depends on the reference group. That is,
	\begin{equation}
        \label{eq:sign_flip_condition_explained_neq}
		\sign\!\left(\Delta\mu^T \beta_H\right) \neq \sign\!\left(\Delta\mu^T \beta_K\right).
	\end{equation}
\end{definition}

From the definition, we can immediately see the following alternative characterization.
%
\begin{remark}
    \label{rmk:explained_ordering_sensitivity}
    The OBD explained component flips signs if and only if the following two conditions hold:
    \begin{equation}
    \label{eq:sign_flip_condition_explained}
        \Delta\mu^T \beta_H \neq \Delta\mu^T \beta_K, \quad \min\{\Delta\mu^T \beta_H,\ \Delta\mu^T \beta_K\} < 0 < \max\{\Delta\mu^T \beta_H,\ \Delta\mu^T \beta_K\}.
    \end{equation}
\end{remark}
Since $\Delta\mu$ is constant across analyses, sign flips in the explained component are straightforwardly driven
by differences in sign in the slope between groups.

% A large region of the parameter space satisfies \cref{eq:sign_flip_condition_explained_neq}.
% We compute the volume of this region by integrating the region that satsfies \cref{eq:sign_flip_condition_explained_neq} against a Uniform(-M, M) measure on \( \beta_H \) and \( \beta_K \), where \( M \) is large.
% The computations, which appear in \cref{app:sign_flip_volume}, show that 50\% of the parameter space satisfies \cref{eq:sign_flip_condition_explained_neq}.

%To see how this characterization applies in practice, we return to the ICU mortality and heart rate example presented in \cref{sec:real_example}. In that example, we observed a sign flip in the explained component of the OBD when using men or women as the reference group. Within HR quartile two, the covariate mean difference $\Delta\mu$ between men and women is fixed by the sample, but the coefficients governing how mortality loads on these covariate differences differ in sign across groups. In particular, the opposite directions of the heart rate coefficients for men and women shown in \cref{fig:hr_quartile2} suggest that a sign flip can occur. As a result, $\Delta\mu^T\beta_{\text{women}}$ and $\Delta\mu^T\beta_{\text{men}}$ have opposite signs, leading to a sign reversal in the explained component. A numerical breakdown of $\Delta\mu$ and the group-specific coefficients is reported in Appendix~\cref{tab:delta_mu_betas}.

\textbf{Unexplained component.}
Unlike the explained component, sign reversals in the unexplained component arise from the interaction between slope and intercept differences.
\begin{definition}
    \label{def:unexplained_sign_flip}
	Recall $\Delta\beta := \beta_H - \beta_K$ and $\Delta\alpha := \alpha_H - \alpha_K$. The OBD \emph{unexplained component flips signs} if the sign of the unexplained term, $\mu_{g'}^T \Delta\beta + \Delta\alpha$, depends on the reference group. That is
	\begin{equation}
		\sign\left( \mu_K^T \Delta\beta + \Delta\alpha \right) \neq \sign\left( \mu_H^T \Delta\beta + \Delta\alpha \right).
	\end{equation}
\end{definition}

The following alternative characterization for sign flips of the unexplained component is analogous to \cref{rmk:explained_ordering_sensitivity}.
%
\begin{restatable}[OBD unexplained sign flips]{proposition}{OBSignFlips}
    \label{thm:unexplained_ordering_sensitivity}
    The OBD unexplained component flips signs if and only if the following two conditions hold:
    \begin{equation}
    \label{eq:sign_flip_condition}
        \mu_H^T \Delta\beta \neq \mu_K^T \Delta\beta, \quad \min\{\mu_H^T \Delta\beta, \mu_K^T \Delta\beta\} < -\Delta\alpha < \max\{\mu_H^T \Delta\beta, \mu_K^T \Delta\beta\}.
    \end{equation}
\end{restatable}    
See \cref{sec:sign_flip_proof} for the proof.
From \cref{thm:unexplained_ordering_sensitivity}, we see that group-specific slopes serve a different role in sign flips in the unexplained case relative to the explained case.
Suppose we hold fixed the covariate means $\mu_H$ and $\mu_K$. Then scaling up the magnitude of the vector change-in-slopes ($\Delta \beta$) induces sign flips for a greater range of values 
of the scalar change-in-intercepts ($\Delta \alpha$).
That is, larger magnitudes of \( \Delta \beta \) create ``more room'' for sign flips in the unexplained component.

%    \Cref{thm:ordering_sensitivity} shows that sign reversals in the unexplained component are not edge cases. Equation \(U^{(K)} - U^{(H)} = -B\) implies that switching the reference group shifts the unexplained term by exactly \(B\), and the explained part by \(-B\). When \(B \neq 0\), this shift can flip the diagnosis—from structural underpayment to overpayment—even with a correctly specified model. The disagreement region \(S \in (-B, 0)\) for \(B > 0\) (and its mirror for \(B < 0\)) is tight: outside this interval, both orderings agree in sign; inside it, they disagree (see \cref{imageTODO}). This window has width \(|B|\), which grows with larger covariate shifts or slope differences—both common in practice. So sign flips are not rare or pathological; they’re a general sensitivity that arises whenever \(B\) is materially different from zero. This motivates the next result.
%
%    The next corollary formalizes that the explained parts across orderings move in perfect opposition to the unexplained parts. Hence, any instability in the structural attribution is mirrored one-for-one by instability in the endowment attribution. This zero-sum coupling is central: it shows that ``credit” assigned to covariate differences depends on the arbitrary choice of reference.
%    
%
%    \begin{restatable}[]{corollary}{explainablepart}
%        \label{cor:explained_flip}
%        With \(B\) as above,
%        \(
%        E^{(K)} = E^{(H)} - B
%        \).
%        Thus the explained components flip sign iff \(E^{(H)}(E^{(H)}-B)<0\); any flip window has width \(|B|\), and
%        \(
%        (E^{(K)}-E^{(H)}) = -(U^{(K)}-U^{(H)})
%        \).
%    \end{restatable}
%
%
%It may be tempting to wonder whether this sensitivity is somehow bounded by the observed outcome gap \(\Delta_Y\). \Cref{cor:unbounded_sensitivity} in the Appendix addresses this question directly: even when \(\Delta_Y\) is held fixed, the ordering-induced swing in the structural component can be made arbitrarily large. That is, no a priori bound ties the magnitude of ordering sensitivity to the magnitude of the empirical gap itself.

\subsection{Are sign flips an adversarial phenomenon?}
\label{sec:sign_flip_fraction}

After characterizing sign flips in \cref{sec:sign_flip_characterizations}, we can ask whether they occur only for adversarially chosen parameters.
One formalization of this question is to ask what proportion of the volume of parameter space corresponds to sign flips --- and in particular, whether that proportion is non-zero.

An immediate challenge with talking about volume within a parameter space is that the parameter space inherits
its units from the covariates and response; in particular, any individual unit choice is somewhat
arbitrary (e.g., centimeters vs.\ millimeters), but also different covariates need not be 
comparable (e.g., temperature vs.\ weight). 
The usual way to alleviate this incompatibility in one group is to center the covariates and set their scale (e.g., standard deviation) to be 1. 
We can use similar normalization to center the covariates such that $\mu_K = \mathbf{0} \in \R^d$ and scale the covariates such that $\Delta\mu = \mathbf{1} \in \R^d$ (i.e., $\mu_H = \mathbf{1}$).
%To alleviate this potential incompatibility between dimensions,
%we first make a common standardization choice in one population: we transform the covariate means to be zero, $\mu_K = \mathbf{0}$, where \( \mathbf{0} \) is a vector of zeros with appropriate length. How to standardize population spread and how to handle the second group are less obvious,
%but we address both of these challenges by choosing $\mu_H = \mathbf{1}$, where \( \mathbf{1} \) is a vector of ones with appropriate length.
Our standardization choice ($\mu_K= \mathbf{0},\mu_H=\mathbf{1}$) is
possible as long as $\mu_K \neq \mu_H$ elementwise (see \cref{lemma:units}).

In the theory that follows, we will find it convenient to restrict our focus to a bounded range of parameters. Namely, we
assume that, for some $M > 0$ and all groups $g$, the entries of $\beta_g$ and $\alpha_g$ lie in $[-M,M]$.
Note, however, that we do not restrict the choice of $M$, and the proportions we compute in all
cases end up having no dependence on $M$.

%To do so, we must choose a region of parameter space we are interested in and then compute the volume of this region that leads to sign flips.
%What region of parameter space should we pick?
%In absence of any specific applied problem, we first standardize units of $X$ such that $\mu_K = 0$ and $\mu_H = 1$ (this is always possible as long as $\mu_K \neq \mu_H$; see \cref{lemma:units}) and then let all entries of $\beta_g$ and $\alpha_g$ lie in $[-M,M]$ for $M > 0$, $g \in \{H,K\}$.
%The fact that each $\mu_g$ has constant entries and that each entry of each $\beta_g$ lies in the same range means that we are making no assumption on the magnitude of each covariate's constribution to the response $Y$ (i.e., $\mu_{g,i}\beta_{g,i}$ sits in the same range for all $i=1,\dots,d$).
As the next two results demonstrate, under these choices, the fraction of parameters leading to sign flips is fairly large.
%
First, for the explained component, we find that, under our conditions, half of the parameter space $(\beta_H, \beta_K, \alpha_H, \alpha_K)$ corresponds to sign flips.
\begin{proposition} \label{prop:explained_fraction}
	Suppose $\mu_K = \mathbf{0}$ and $\mu_H = \mathbf{1}$.
	Choose any $M > 0$.
	Let $C_M \subset \R^{2d}$ be a cube with side length $2M > 0$ centered at the origin, and suppose $(\beta_H, \beta_K)$ jointly lie in this cube. Let $C_{flip} \subset C_M$ be the settings of $(\beta_H, \beta_K)$ that lead to a sign flip in the explained component.
	Then
	\begin{equation}
		\frac{\mathrm{Volume}(C_{flip})}{\mathrm{Volume}(C_M)} = \frac{1}{2}.
	\end{equation}
\end{proposition}
See \cref{proof:explained_fraction} for the proof.

Second, for the unexplained component, we find that the proportion of parameter space corresponding to sign flips nears 1/2 as the dimension of the covariate vector grows.
\begin{proposition} 
\label{prop:unexplained_fraction}
	Suppose $\mu_K = \mathbf{0}$ and $\mu_H = \mathbf{1}$.
	Let $C_M \subset \R^{2d+2}$ be a cube with side length $2M > 0$ centered at the origin, and suppose $(\beta_K, \beta_H, \alpha_K, \alpha_H)$ jointly lie in this cube.
	Let \( I_2 \sim \text{Irwin-Hall}(2) \) and \( J_{2d} \sim \text{Irwin-Hall}(2d) \) be independent random variables.
	The fraction of $C_M$ for which a sign flip occurs in the unexplained component is given by:
\begin{align} \label{eq:unexplained_fraction}
    P_d := \mathrm{Pr} \left( \{ I_2 > 1 \} \cap \{ J_{2d} < d + 1 - I_2 \}  \right) + \mathrm{Pr} \left( \{ I_2 < 1 \} \cap \{ J_{2d} > d + 1 - I_2 \} \right),
\end{align}
where the probability $\mathrm{Pr}$ is over the independent random variables $I_2, J_{2d}$.
Moreover, $P_d \underset{d \to \infty}{\longrightarrow} 1/2$.
\end{proposition}
See \cref{proof:unexplained_fraction} for the proof.
The joint Irwin-Hall probability in \cref{eq:unexplained_fraction} is hard to evaluate analytically, so we evaluate it numerically for various values of $d$; see \cref{fig:sign_flip_overview} (right).
For even modest dimension $d$, the fraction of parameter space leading to sign flips exceeds 40\%.
% Figure moved to \cref{fig:sign_flip_overview} (right) in \cref{sec:real_example}.
% \begin{figure}
%     \centerline{\includegraphics[width=0.75\linewidth]{Sections/Figures/prob_unexplained.pdf}}
%     \caption{Percentage of parameter space leading to sign flips. See \cref{remark:computing_irwin_hall} for how we compute the percentage for the unexplained component.}
%     \label{fig:prob_unexplained}
% \end{figure}

While our results above assume the standardization $\mu_K=\mathbf{0}, \mu_H=\mathbf{1}$, 
we confirm by simulations that we expect similar results when we take each coordinate of $\mu_g$ uniform
in the range $[-M,M]$ as well. See \cref{appendix:sign_flip_simulation} for details.

\subsection{Reconciling the theoretical and empirical prevalence of sign flips}
\label{sec:reconciling_theory_empirics}

Our main focus in \cref{sec:sign_flip_fraction} was establishing that a non-zero proportion of parameter
space yielded sign flips.
% In fact, our theoretical results suggest a large fraction of parameter space produces sign flips in both OBD components.
One limitation of this theory is that it assumes regression parameters are uniformly distributed, when in practice we may expect they follow non-uniform distributions.
Indeed, our earlier empirical exercise from \cref{tab:flip_summary_main}, finds that sign flips are less common than a direct application of theory would suggest. 
Even if we count total flips and ignore statistical significance, we see sign flips in 19\% of (i.e., 27 of 139) tested subsets for the OBD.
We now describe and interpret constraints on the parameter space that may produce fewer sign flips than implied by our theoretical results.

% An exhaustive search of the literature using the OBD is beyond the scope of the present paper;
% a paper corresponding to a Stata implementation of the OBD has thousands of citations alone \citep{jann2008}, and there are implementations
% in other languages as well \citep{Hlavac2022oaxaca}.
% Without trying to form a representative sample, we examined a selection of papers performing the OBD. 
% Of the papers we examined, many did not report both directions of the OBD.
% Of the papers that reported both directions \citep{ONeill2006, fortinLemieuxFirpo2011, sharafRashad2016, rahimiNazari2021}, none exhibited a sign flip.

%% AS: this is now introduced (briefly) in Sec 3
% Publication or other biases could affect the published literature. To avoid these concerns and also offer a more systematic investigation,
% we examined collections
% of OBD analyses one might potentially run on two publicly available datasets. 
% In each case, our systematic search for sign flips turned up few examples.
% First, we examined the PhysioNet ICU cohort collected for the study of in-hospital mortality \citep{goldberger2000physiobank}.
% We started by evaluating the OBD across nine predefined data subsets: the full cohort, each ICU type category, two age groups (patients younger than $65$ and patients aged $65$ or older), and two groups defined by a median split of the SAPS--I score. For each of these subsets, we fit gender-specific linear models and computed the OBD. 
% We found no sign flips.
% We then expanded our analysis to a larger and more heterogeneous collection of subpopulations.
% Across all $139$ data subsets, we found that $27$ exhibited at least one sign reversal --- a far cry from the number predicted by \cref{prop:explained_fraction,prop:unexplained_fraction}.
% The HR quartile two example in \cref{sec:real_example} was one such example. We describe the data and our search procedure fully in 
% \cref{appendix:icu_details}.

% We ran a similar search on U.S.\ labor force data; in particular,
% the raw data come from the 2016 American Community Survey 1-Year Public Use Microdata Sample \cite{uscensus_acs_pums_2016}.
% We again purposely constructed a large set of data analyses to search for sign flips in different subsets of the data, comprising all combinations of the factors in Appendix \cref{tab:us_census_design}.
% Out of the 355 comparisons we examined, we found zero cases where the explained component flipped sign and one case where the unexplained component flipped sign.
% See \cref{appendix:us_labor_force} for details of this data set and our search procedure.
% Again, these numbers are substantially lower than those predicted by \cref{prop:explained_fraction,prop:unexplained_fraction}.

% We conclude from these two searches that we do not expect sign flips to be more than empirically rare, but
% their existence is still possible.

% How can it be that sign flips appear to be rare in practice (\cref{sec:sign_flips_practice}) but common in terms of parameter-space volume (\cref{sec:sign_flip_fraction})?

\textbf{Explained component.} First, consider the explained component. Under the standardization from \cref{sec:sign_flip_fraction},
we can simplify the characterization of when a sign flip happens.
%
\begin{remark} \label{rmk:no_explained_sign_flip}
	Suppose $\mu_K = \mathbf{0}$ and $\mu_H = \mathbf{1}$.
	Then there is a sign flip if and only if $\sign (\mathbf{1}^\top \beta_{H}) \neq \sign( \mathbf{1}^\top \beta_{K})$.
\end{remark}
%
This observation follows immediately from \cref{defn:sign_flip_explained}.
In some contexts, we may expect the elements of $\beta_g$ to be of the same sign in both groups.
Consider the ICU mortality example from \cref{sec:real_example}.
While we may expect risk factors to have different effect sizes for men and women a priori, there is little reason to think that many factors that increase men's mortality would decrease women's (or vice versa).
\cref{rmk:no_explained_sign_flip} could then explain why we observe explained component sign flips in only 11 out of the 139 subgroups we test (\cref{tab:flip_summary_explained}).
% AS: the flips could also reflect noise — the subgroups we test might be small, so that they _don't_ present much evidence about population-level differences (in line with few statistically significant flips)
% Indeed, this observation makes the explained component sign flip in \cref{sec:real_example} more surprising, as it implies substantive differences in mortality risk in this subpopulation of men and women.

\textbf{Unexplained component.} Next, consider the unexplained component. In this case, a sign flip intuitively requires the intercept and slope effects to pull in opposite directions across groups. When both move in the same direction, so that the group with higher baseline outcomes also exhibits higher returns to covariates, the unexplained component preserves its sign regardless of the chosen reference group.
We now formalize these ideas.
\begin{assumption}[Aligned intercept--slope effects]
\label{ass:no_sign_flip}
%\[
%    (\mu_H^T \Delta\beta)(\mu_K^T \Delta\beta) > 0 
%\quad \text{and}  
%\]
$
    \sign(\Delta\alpha) = \sign(\mu_H^T \Delta\beta) = \sign(\mu_K^T \Delta\beta)
$
\end{assumption}

\begin{proposition}
    \label{prop:no_sign_flip}
    Under \cref{ass:no_sign_flip}, the unexplained components 
    \( U^{(H)}\) and  \(U^{(K)}\) share the same sign.
\end{proposition}
See \cref{proof:no_sign_flip} for the proof.
In some domains, intercept and slope advantages might often move in the same direction; 
\cref{prop:no_sign_flip} rules out sign flips in such cases.
To illustrate, consider an analysis of the gender wage gap that attributes the difference in average earnings between men and women to covariates (e.g., work experience) and to outcomes given covariates (e.g., how experience is remunerated).
\cref{ass:no_sign_flip} corresponds to the group with higher baseline wages ($\alpha_g$) also having higher returns to experience ($\beta_g$). 
\cref{prop:no_sign_flip} is one explanation for why sign flips are rare in our analysis of U.S. labor force data (\cref{appendix:us_labor_force}): \cref{ass:no_sign_flip} holds in 182 out of 400 comparisons we search (45.5\%), ruling out unexplained component sign flips in those cases.

However, \cref{ass:no_sign_flip} is not plausible in all domains.
In the ICU setting from \cref{sec:real_example}, the assumption requires the higher mortality population to also be more sensitive to observed covariates.
This relationship may fail to hold when important factors --- such as environmental exposure, chronic conditions, or family history of disease --- are not recorded in a researcher's dataset.
If such \textit{omitted covariates} are strong predictors of mortality, a population may have a high mortality rate that is less sensitive to \textit{observed covariates}.
Consistent with this intuition, \cref{ass:no_sign_flip} never holds across the 139 comparisons in the ICU data exercise.
In \cref{appendix:ovb}, we explore how the presence of omitted covariates can generate parameter differences $\Delta\beta$ and $\Delta\alpha$ that do preclude sign flips, without invoking \cref{ass:no_sign_flip}.
%We develop this argument with more rigor in \cref{appendix:ovb}, explicitly showing how the regression parameters \( \alpha \) and \( \beta \) are jointly influenced by unobserved covariates.

% On the other hand, sign flips remain plausible for the unexplained component. 
% Consider the SBP--BMI example in \cref{sec:simulated_example}, \cref{ass:no_sign_flip} would correspond to a case where the higher-resource group exhibits both a higher baseline SBP and a stronger SBP--BMI relationship than the lower-resource group. However, in the simulated data, the lower-resource group shows a steeper SBP--BMI slope while having a lower baseline; that is, the SBP--BMI example violates \cref{ass:no_sign_flip}, and indeed we see a sign flip in the unexplained component.

\subsection{Proof of \cref{thm:unexplained_ordering_sensitivity}}
\label{sec:sign_flip_proof}
\begin{proof}
Using \cref{def:unexplained_sign_flip}, we want to show that 
\[
\sign\!\left( \mu_K^T \Delta\beta + \Delta\alpha \right)
\neq
\sign\!\left( \mu_H^T \Delta\beta + \Delta\alpha \right)
\iff
\begin{aligned}
&\mu_H^T \Delta\beta \neq \mu_K^T \Delta\beta, \\
&\min\{\mu_H^T \Delta\beta, \mu_K^T \Delta\beta\}
<
-\Delta\alpha
<
\max\{\mu_H^T \Delta\beta, \mu_K^T \Delta\beta\}.
\end{aligned}
\]

$\Rightarrow$: Suppose that $\sign\left( \mu_K^T \Delta\beta + \Delta\alpha \right) \neq \sign\left( \mu_H^T \Delta\beta + \Delta\alpha \right)$. If $\mu_H^T \Delta\beta = \mu_K^T \Delta\beta$, then both $\mu_g^T \Delta\beta + \Delta\alpha$ would have the same sign, which contradicts our assumption of opposite signs. Thus, we must have $\mu_H^T \Delta\beta \neq \mu_K^T \Delta\beta$.

Next, suppose that $\mu_K^T \Delta\beta + \Delta\alpha >0$ and $\mu_H^T \Delta\beta + \Delta\alpha <0$. Then, 
\[
\mu_K^T \Delta\beta > -\Delta\alpha \quad \text{and} \quad \mu_H^T \Delta\beta < -\Delta\alpha.
\]
Combining these two inequalities, we have 
\begin{equation}
  \label{eq:intermediate_ineq1}
  \mu_H^T \Delta\beta < -\Delta\alpha < \mu_K^T \Delta\beta.
\end{equation}
Similarly, if $\mu_K^T \Delta\beta + \Delta\alpha <0$ and $\mu_H^T \Delta\beta + \Delta\alpha >0$, then
\begin{equation}
  \label{eq:intermediate_ineq2}
  \mu_K^T \Delta\beta < -\Delta\alpha < \mu_H^T \Delta\beta.
\end{equation}
Combining \cref{eq:intermediate_ineq1,eq:intermediate_ineq2}, we conclude that
\[
  \min\{\mu_H^T \Delta\beta, \mu_K^T \Delta\beta\} < -\Delta\alpha < \max\{\mu_H^T \Delta\beta, \mu_K^T \Delta\beta\}.
\]

$\Leftarrow$; We now show the converse direction. Suppose that:

$(1)$ $\mu_H^T \Delta\beta \neq \mu_K^T \Delta\beta$ and 

$(2)$ $\min\{\mu_H^T \Delta\beta, \mu_K^T \Delta\beta\} < -\Delta\alpha < \max\{\mu_H^T \Delta\beta, \mu_K^T \Delta\beta\}$. 

We want to show that $\sign\left( \mu_K^T \Delta\beta + \Delta\alpha \right) \neq \sign\left( \mu_H^T \Delta\beta + \Delta\alpha \right)$.
$(2)$ implies that one of the following must be true:
\[
  \mu_H^T \Delta\beta < -\Delta\alpha < \mu_K^T \Delta\beta, \quad \text{or} \quad \mu_K^T \Delta\beta < -\Delta\alpha < \mu_H^T \Delta\beta.
\]
If the first one holds, then $\mu_H^T \Delta\beta + \Delta\alpha <0$ and $\mu_K^T \Delta\beta + \Delta\alpha >0$, that is, the unexplained components have opposite signs. If the second holds, a symmetric argument yields the same conclusion as desired. 

% First, consider the case where $\sign \mu_K^T \Delta\beta = \sign \mu_H^T \Delta\beta$.
% In this case, the addition of $\Delta\alpha$ must be small enough to flip one of the signs of the $\mu_g^T \Delta\beta$ but so large that it flips the other.
% As we have assumed $\mu_K^T \Delta\beta \neq \mu_H^T \Delta\beta$, one of the $\mu_g^T \Delta\beta$ will be lesser in magnitude; without loss of generality, suppose this is $g = H$.
% That is, $\abs{\mu_H^T \Delta\beta} < \abs{\mu_K^T \Delta\beta}$.
% Now, we have also assumed that $\sign \Delta\alpha \neq \sign \mu_K^T \Delta\beta$ and $\abs{\mu_H^T \Delta\beta} < \abs{\Delta\alpha} < \abs{\mu_K^T \Delta\beta}$.
% We conclude that $\sign( \mu_H^T \Delta\beta + \Delta\alpha) \neq \sign( \mu_K^T \Delta\beta + \Delta\alpha)$.

% Next, consider the case where $\sign \mu_K^T \Delta\beta \neq \sign \mu_H^T \Delta\beta$.
% In this case, the addition of $\Delta\alpha$ to the $\mu_g^T \Delta\beta$ must be small enough such that it does not change the sign of either one.
% Without loss of generality, suppose that $\sign \Delta\alpha \neq \sign \mu_H^T \Delta\beta$.
% Then the assumed condition $\abs{\mu_H^T \Delta\beta} > \abs{\Delta\alpha}$ is enough to ensure that $\sign(\mu_H^T \Delta\beta) = \sign(\mu_H^T \Delta\beta + \Delta\alpha) \neq \sign(\mu_K^T \Delta\beta + \Delta\alpha)$.
\end{proof}

% ---------------------------------------------------------
% \mqedit{
% \subsection{Decision tree for sign flips}
% \label{sec:decision_tree}

% To shed more light on \cref{thm:ordering_sensitivity}, we disentangle the condition in \cref{eq:sign_flip_condition} into a decision tree of when a sign flip can occur: 

% \begin{algorithm}[H]
%     \caption{Unexplained component sign-flip}
%     \label{alg:decision_tree}
% \begin{algorithmic}  
%     \If{$\mu_H^T \Delta\beta \neq \mu_K^T \Delta\beta$}
%         \If{$\sign(\mu_H^T \Delta\beta) = \sign(\mu_K^T \Delta\beta)$}
%             \If{$\sign(\Delta\alpha) \neq \sign(\mu_H^T \Delta\beta)$}
%                 \If{$|\mu_H^T \Delta\beta| < |\Delta\alpha| < |\mu_K^T \Delta\beta|$ \textbf{or} $|\mu_K^T \Delta\beta| < |\Delta\alpha| < |\mu_H^T \Delta\beta|$}
%                     \State A sign flip occurs.
%                 \EndIf
%             \EndIf
%         \ElsIf{$\sign(\mu_H^T \Delta\beta) \neq \sign(\mu_K^T \Delta\beta)$}
%             \If{$\sign(\Delta\alpha) = \sign(\mu_H^T \Delta\beta)$ \textbf{and} $|\mu_K^T \Delta\beta| > |\Delta\alpha|$}
%                 \State A sign flip occurs.
%             \ElsIf{$\sign(\Delta\alpha) \neq \sign(\mu_H^T \Delta\beta)$ \textbf{and} $|\mu_H^T \Delta\beta| > |\Delta\alpha|$}
%                 \State A sign flip occurs.
%             \EndIf
%         \EndIf
%     \Else
%         \State There is no sign flip.
%     \EndIf
% \end{algorithmic}
% \end{algorithm}

% }

\subsection{Unbounded sensitivity at fixed mean gap}
 \label{ex:unbounded}
In this section we present a toy example where we can have arbitrarily large flip magnitudes. Consider the following example in dimension \( 1 \), where we assume we have a parameter \( L > 0 \) which controls the magnitude of the covariate shift, and \( \varepsilon > 0 \) which controls the deviation in slope coefficients. We assume the mean vector is given by \( (\mu_K, \mu_H) = (0, L) \), and the slope parameters deviate symmetrically from \( \frac{1}{L} \), so that \( \beta_H = \frac{1}{L} + \frac{\varepsilon}{2} \), \( \beta_K = \frac{1}{L} - \frac{\varepsilon}{2} \). Then, the explained and unexplained components under both reference groups are:
\begin{align*}
    E^{(H)} &= \Delta\mu \cdot \beta_H = L\left(\frac{1}{L} + \frac{\varepsilon}{2}\right) = 1 + \frac{L\varepsilon}{2}, \quad U^{(H)} = \mu_K \cdot \Delta\beta = -\frac{L\varepsilon}{2}, \\
    E^{(K)} &= \Delta\mu \cdot \beta_K = L\left(\frac{1}{L} - \frac{\varepsilon}{2}\right) = 1 - \frac{L\varepsilon}{2}, \quad U^{(K)} = \mu_H \cdot \Delta\beta = +\frac{L\varepsilon}{2}.
\end{align*}

Hence, the difference between the two decompositions is \( E^{(H)} - E^{(K)} = L\varepsilon \), which can be made arbitrarily large as \( L \to \infty \), even though the total mean gap is always fixed at \( \Delta_Y = 1 \). Similarly, the unexplained components differ by \( U^{(K)} - U^{(H)} = L\varepsilon \).

\begin{center}
\begin{tabular}{ccccccc}
\toprule
\(L\) & \(\varepsilon\) 
& \(E^{(H)}\) & \(E^{(K)}\) 
& \(U^{(H)}\) & \(U^{(K)}\) 
& Flip? \\
\midrule
20    & 0.1  & \(2\)     & \(0\)      & \(-1\)    & \(1\)     & \checkmark \\
200   & 0.1  & \(11\)    & \(-9\)     & \(-10\)   & \(10\)    & \checkmark \\
1000  & 0.1  & \(51\)    & \(-49\)    & \(-50\)   & \(50\)    & \checkmark \\
\bottomrule
\end{tabular}
\end{center}

Note that the sign flip is symmetric by construction, but the magnitude of the discrepancy grows linearly in \( L \), highlighting that the attribution gap between reference groups can be stretched arbitrarily far.

\subsection{Decision Tree Representation of \cref{eq:sign_flip_condition}}

To illustrate \cref{thm:unexplained_ordering_sensitivity}, we represent the condition in \cref{eq:sign_flip_condition} as a decision tree in \cref{alg:decision_tree} and use it to scrutinize the healthcare example in \cref{sec:simulated_example}.
In this example, $\mu_H$ and $\mu_K$ denote mean BMIs, both of which are positive by definition and thus $\sign{(\mu_H \Delta\beta)} = \sign{(\mu_K \Delta\beta)} = \sign{(\Delta\alpha)}$. 
Therefore, following \Cref{alg:decision_tree}, a sign flip will occur only if $\sign{(\Delta\alpha)} \neq \sign{(\mu_H \Delta\beta)} = \sign{(\Delta\beta)}$; equivalently, it must hold that one group has a larger $\beta_g$ (more sensitive to higher BMI) but lower $\alpha_g$ (lower SBP on average).
This might occur if one group has better access to treatment but worse genetic factors, as in \cref{sec:simulated_example}.
Furthermore, a sign flip will occur if either $\abs{\mu_H \Delta\beta} < \abs{\Delta\alpha} < \abs{\mu_K \Delta\beta}$ or $\abs{\mu_K \Delta\beta} < \abs{\Delta\alpha} < \abs{\mu_H \Delta\beta}$. This condition implies a range of values over which sign flips will occur; that is, there will be at least one choice of reference group under which improved healthcare appears to help in lowering mean SBP and another ordering under which improved healthcare appears to \emph{hurt} mean SBP.

Even though \cref{thm:unexplained_ordering_sensitivity} does not restrict the magnitude of the parts of the unexplained components, in \cref{ex:unbounded} we provide an example to illustrate that these sign flips are not relegated to settings in which $\mu_g^T\Delta\beta$ are small in magnitude; in fact, we show that they can be arbitrarily large in magnitude.

  \begin{algorithm}
    \caption{Unexplained component sign flip}
    \label{alg:decision_tree}
    \begin{algorithmic}  
        \If{$\mu_H^T \Delta\beta \neq \mu_K^T \Delta\beta$}
            \If{$\sign(\mu_H^T \Delta\beta) = \sign(\mu_K^T \Delta\beta)$}
                \If{$\sign(\Delta\alpha) \neq \sign(\mu_H^T \Delta\beta)$}
                    \If{$|\mu_H^T \Delta\beta| < |\Delta\alpha| < |\mu_K^T \Delta\beta|$ \textbf{or} $|\mu_K^T \Delta\beta| < |\Delta\alpha| < |\mu_H^T \Delta\beta|$}
                        \State A sign flip occurs.
                    \Else
                    	\State There is no sign flip.
                    \EndIf
                \Else
                	\State There is no sign flip.
                \EndIf
            \ElsIf{$\sign(\mu_H^T \Delta\beta) \neq \sign(\mu_K^T \Delta\beta)$}
                \If{$\sign(\Delta\alpha) = \sign(\mu_H^T \Delta\beta)$ \textbf{and} $|\mu_K^T \Delta\beta| > |\Delta\alpha|$}
                    \State A sign flip occurs.
                \ElsIf{$\sign(\Delta\alpha) \neq \sign(\mu_H^T \Delta\beta)$ \textbf{and} $|\mu_H^T \Delta\beta| > |\Delta\alpha|$}
                    \State A sign flip occurs.
                \Else
                    \State There is no sign flip.
                \EndIf
            \EndIf
        \Else
            \State There is no sign flip.
        \EndIf
    \end{algorithmic}
\end{algorithm}

\subsection{Relationship Between Unexplained Component Sign Flips and Omitted Variables}
\label{appendix:ovb}
Here, we assume the same standardization of covariates as in \cref{sec:sign_flip_fraction}; that is, we assume $\mu_K = \mathbf{0}$ and $\mu_H = \mathbf{1}$.
Under this normalization, unexplained component sign flips require \( | \mathbf{1}^T \Delta \beta | > | \Delta \alpha | \), so any data-generating processes that jointly produces large \( | \mathbf{1}^T \Delta \beta | \) and \( | \Delta \alpha | \) would make sign flips less likely than \cref{prop:unexplained_fraction} suggests.
In this section, we present a mechanism that can jointly produce large \( | \mathbf{1}^T \Delta \beta | \) and \( | \Delta \alpha | \) without invoking \cref{ass:no_sign_flip}.

The mechanism involves the researcher's dataset omitting covariates that are related to both \( X \) and \( Y \).
For example, hospital records only contain measurements that doctors have ordered, not the results of important tests a doctor may have skipped.
Similarly, census data used for studying wage gaps typically lack detailed information about skills and experience from a worker's resume and work history.
Finally, in practice, nonlinear transformations of the covariates are also effectively omitted from analysis, like when a researcher does not include appropriate polynomial terms in a regression specification.
Let \( Z \in \mathbb{R}^{d'} \) be a vector of such \textit{omitted covariates}.
We will establish that both \( \Delta \beta \) and \( \Delta \alpha \) can depend on the relationship between \( Z, X, \) and \( Y \).

Our argument requires us to specify the conditional expectations \( \E_g[Y \mid X, Z] \) and \( \E_g[Z \mid X] \).
To establish a simple example, we will assume that both are linear and satisfy
\begin{align}
  \label{eq:y_on_x_z}
  \E_g[Y \mid X, Z] &= \omega_g + X^\top \theta_g + Z^\top \gamma_g, \\
  \label{eq:z_on_x}
  \E_g[Z \mid X] &= \zeta_g + \Psi_g X,
\end{align}
where $\omega_g \in \R$, $\theta_g \in \R^d$, $\gamma_g \in \R^{d'}$, $\zeta_g \in \R^{d'}$, and $\Psi_g \in \R^{d' \times d}$.
Relaxing these assumptions may lead to other non-uniform relationships between \( \Delta \beta \) and \( \Delta \alpha \).

Recall that the OBD is based on regressing the outcome against the observed covariates for each group \(g \in \{H, K\}\) via
\begin{align}
    \label{eq:y_on_x}    
    \E_g[Y \mid X] = \alpha_g + X^\top \beta_g.
\end{align}
We will start by showing that \( \alpha_g \) and \( \beta_g \) depend on \( \gamma_g, \zeta_g, \), and \( \Psi_g \). 
To do so, apply the law of iterated expectations to \cref{eq:y_on_x}, plug in \cref{eq:y_on_x_z,eq:z_on_x}, and then rearrange terms:
\begin{align}    
    \E_g[Y \mid X] 
    &= \E_g[ \E_g[Y \mid X, Z] \mid X ] \nonumber \\
    &= \E_g[ \omega_g + X^\top \theta_g + Z^\top \gamma_g \mid X ] \nonumber  \\
    &= \omega_g + X^\top \theta_g + \E_g[Z \mid X]^\top \gamma_g \nonumber \\
    &= \omega_g + X^\top \theta_g + (\zeta_g + \Psi_g X)^\top \gamma_g \nonumber \\ 
    \label{eq:ovb}
    &= \underbrace{\omega_g + \zeta_g^\top \gamma_g}_{=\alpha_g} + \E_g[X]^\top \underbrace{(\theta_g + \Psi_g^\top \gamma_g)}_{=\beta_g}.
\end{align}
\cref{eq:ovb} is the well-known \textit{omitted variables bias} formula.
It describes how the slope and intercept of a regression equation depend on variables that are not included in the equation.
We can use this formula to express the difference in slopes and intercepts across the two populations as 
\begin{align}  
  \label{eq:delta_alpha_gamma}
  \Delta \alpha &= \Delta \omega + \Delta \zeta^\top \Delta \gamma,  \\
  \label{eq:delta_beta_gamma}
  \Delta \beta  &= \Delta \theta + \Delta \Psi^\top \Delta \gamma,
\end{align}
where \( \Delta \) takes the $H-K$ difference in the quantity that follows; for example, \( \Delta \omega := \omega_H - \omega_K \).
These equations show that \( \Delta \alpha \) and \( \Delta \beta \) depend on the role of omitted covariates in the data-generating process.
Crucially, they provide a mechanism to produce \( | \Delta \alpha | \) and \( | \mathbf{1}^\top \Delta \beta | \) that are jointly large:
if both the Z-X and Z-Y relationships change in complementary ways, then both \( \Delta \zeta^\top \Delta \gamma \) and \( \mathbf{1}^\top \Delta \Psi^\top \Delta \gamma \) will be large. 
If these terms are not counterbalanced by \( \Delta \omega \) and \( \Delta \theta \), the result will be large \( | \Delta \alpha | \) and \( | \mathbf{1}^\top \Delta \beta | \).

\subsection{Volume of the Sign-Flip Parameter Region}
We first show how to choose units for the covariates conveniently, so that \( \Delta \mu = \mathbf{1} \).
This will simplify volume calculations for the regions of parameter space that flip the explained and unexplained components of the OBD.
\begin{lemma}
  \label{lemma:units}
  If \( \mu_H \neq \mu_K \) elementwise, we can choose units for \( X \) such that \( \mu_K := \E_K[X] = \mathbf{0} \) and \( \mu_H := \E_H[X] = \mathbf{1} \).
  \begin{proof}  
  Consider the transformation \( \tilde{X} = (X - \mu_K) / (\mu_H - \mu_K) \), where the division is elementwise.
  This is well-defined as long as \( \mu_H \neq \mu_K \) elementwise.
  Since \( \tilde{X} \) is a linear transformation of \( X \), it corresponds to a change in units.
  Under these new units, \( \E_K[\tilde{X}] = \mathbf{0} \) and \( \E_H[\tilde{X}] = \mathbf{1} \).  
  \end{proof}
\end{lemma}

\subsubsection{Proof of \cref{prop:explained_fraction}}
\label{proof:explained_fraction}
Apply \cref{lemma:units} so that \( \E_K[X] = \mathbf{0} \), \( \E_H[X] = \mathbf{1} \), and therefore \( \Delta \mu = \mathbf{1} \).
\cref{eq:sign_flip_condition_explained_neq} now simplifies to 
\[ \sign\!\left(\mathbf{1}^\top \beta_{H}\right) \neq \sign\!\left(\mathbf{1}^\top \beta_{K}\right). \]
A Uniform$(-M, M)$ distribution on \( \beta_H \) and \( \beta_K \) corresponds to distributions on \( \mathbf{1}^\top \beta_{H} \) and \( \mathbf{1}^\top \beta_{K} \) that are identical and symmetric around zero.
As such, the signs of \( \mathbf{1}^\top \beta_{H} \) and \( \mathbf{1}^\top \beta_{K} \) will disagree in 50\% of the parameter space.

\subsubsection{Proof of \cref{prop:unexplained_fraction}}
\label{proof:unexplained_fraction}
Apply \cref{lemma:units} so that \( \E_K[X] = \mathbf{0} \), \( \E_H[X] = \mathbf{1} \), and therefore \( \Delta \mu = \mathbf{1} \).
For the rescaled covariates, the conditions in \cref{thm:unexplained_ordering_sensitivity} simplify.
In particular, the unexplained component flips sign if and only if \( \mathbf{1}^\top \Delta \beta \neq 0 \), \( \sign (\Delta \alpha) \neq \sign (\mathbf{1}^\top \Delta \beta) \), and \( | \mathbf{1}^\top \Delta \beta | > |\Delta \alpha| \).

To compute the fraction of $C_M$ for which there are sign flips, we equivalently compute the probability of a sign flip under an independent $\mathrm{Uniform}(-M,M)$ measure on each element of the vector $(\beta_H, \beta_K, \alpha_H, \alpha_K) \in \R^{2d+2}$.
The condition \( \mathbf{1}^\top \Delta \beta \neq 0 \) holds almost surely under the assumed Uniform$(-M,M)$ measure, since \( \mathbf{1}^\top \Delta \beta \) is a non-degenerate linear combination of continuously distributed random variables, so we can focus on the latter two conditions.
To do so, we will determine the distributions of  \( \Delta \alpha \) and \( \mathbf{1}^\top \Delta \beta \).

The sum of \( n \) Uniform$(0,1)$ random variables is sometimes called the Irwin-Hall(\(n\)) distribution.%
\footnote{For \(n = 2\) this is a familiar triangular distribution.}
When we place a Uniform$(-M, M)$ measure on the \( \alpha \)'s and \( \beta \)'s, there is a straightforward relationship between \( \Delta \alpha \), \( \mathbf{1}^\top \Delta \beta \), and appropriately defined Irwin-Hall random variables.
The difference \( \Delta \alpha \) is the sum of two Uniform$(-M, M)$ variables, which is equivalent to an Irwin-Hall($2$) distribution, rescaled from $[0,2]$ to $[-2M, 2M]$.
Similarly, \( \mathbf{1}^\top \Delta \beta \) is the sum of \( 2 d \) Uniform$(-M, M)$ variables, and its distribution is an Irwin-Hall($2 d$) distribution, rescaled from $[0,2d]$ to \( [-2dM, 2dM] \).

That is, let \( I_2 \sim \text{Irwin-Hall}(2) \) and \( J_{2d} \sim \text{Irwin-Hall}(2d) \) be independent.
We can represent
\begin{align*}
    &\mathbf{1}^\top\Delta \alpha = 2M (I_2 - 1) 
    &\mathbf{1}^\top\Delta \beta  = 2M (J_{2d} - d)
\end{align*}

We can use this correspondence to translate the event 
\[
\{ |\mathbf{1}^\top \Delta \beta| > |\Delta \alpha| \}
\cap
\{ \sign(\Delta \alpha) \neq \sign(\mathbf{1}^\top \Delta \beta) \}.
\]
into an equivalent event in terms independent Irwin-Hall random variables:
\begin{align*}
    Pr& \left( \{ |\mathbf{1}^\top \Delta \beta| > |\Delta \alpha| \}
    \cap
    \{ \sign(\Delta \alpha) \neq \sign(\mathbf{1}^\top \Delta \beta) \} \right) \\
    =& 
    Pr \left( \{ \Delta \alpha > 0 \} \cap \{ \mathbf{1}^\top \Delta \beta < - \Delta \alpha \} \right)
    +
    Pr \left(
    \{ \Delta \alpha < 0 \} \cap \{ \mathbf{1}^\top\Delta \beta > - \Delta \alpha \} \right) \\
    =& 
    Pr \left( \{ 2M (I_2 - 1)  > 0 \} \cap \{ 2M (J_{2d} - d) < - 2M (I_2 - 1)  \} \right)  \\
    &+
    Pr \left( \{ 2M (I_2 - 1)  < 0 \} \cap \{ 2M (J_{2d} - d) > - 2M (I_2 - 1)  \} \right) \\
    =& 
    Pr \left( \{ I_2 > 1 \} \cap \{ J_{2d} < d + 1 - I_2 \}  \right) +
    Pr \left( \{ I_2 < 1 \} \cap \{ J_{2d} > d + 1 - I_2 \} \right).
\end{align*}

That is, unexplained component sign flips have the same volume as the region where \( I_2 \) falls above its median and \( J_{2d} \) falls sufficiently far in its left tail, or vice versa.

Finally, as \( d \to \infty \), sign flips have a volume equal to
\begin{equation}
\label{eq:irwin_hall_pr}
 Pr \left( \{ I_2 > 1 \} \cap \{ J_{2d} < d \}  \right) +
    Pr \left( \{ I_2 < 1 \} \cap \{ J_{2d} > d \} \right).
\end{equation}
Since the median of \( I_2 \) is $1$ and the median of \( J_{2d} \) is $d$, and the two random variables are independent, each probability equals \( 0.25 \).
The corresponding fraction of $C_M$ leading to sign flips is therefore \( 0.5 \) as claimed.

\begin{remark} \label{remark:computing_irwin_hall}
We can calculate the probabilities in \cref{eq:irwin_hall_pr} exactly for small \( d \) by using the Irwin-Hall$(n)$ cumulative distribution function (CDF),
\begin{equation}
F(x) = \frac{1}{n!} \sum_{k=0}^{\lfloor x \rfloor} (-1)^k \binom{n}{k} (x - k)^n. \label{eq:exact_irwin_hall}
\end{equation}
However, for larger \( d \), the combinatorial terms in the CDF become difficult to evaluate numerically.
In those cases the central limit theorem allows us to approximate an Irwin-Hall$(2d)$ distribution with a \( N(d, \frac{d}{6})\) distribution.
We compute the probabilities in \cref{sec:sign_flip_fraction} for $d \leq 40$ using the exact CDF in \cref{eq:exact_irwin_hall}, and we use the normal approximation for $d > 40$.
\end{remark}

\subsubsection{Volume of Sign Flip Region Without Covariate Standardization}
\label{appendix:sign_flip_simulation}
The preceding analysis calculates the volume of the sign-flip region when covariates are standardized via \cref{lemma:units}.
Here, we repeat the exercise without standardizing any of the covariates.
Let each of \( \alpha_H, \alpha_K, \beta_H, \beta_K, \mu_H \) and \( \mu_K \) lie inside a cube of appropriate dimension, with side length \( 2M > 0 \) and center at the origin.
Deriving a closed form expression for the volume of the sign flip region is challenging, so we estimate this volume by simulation.
For each \( d \in \{1, \dots, 100 \} \) and \( 2M \in \{2, 20, 200\} \), we draw each regression parameter independently from an appropriate Uniform distribution, and then determine whether the unexplained component flips signs.
We simulate 1,000 draws for each combination of \( d \) and \( M \), and plot the proportion of unexplained component sign flips for each combination in \cref{fig:prob_unexplained_simulated}.

For the explained component, the percentage of the parameter space with sign flips appears constant at 50 \%, regardless of \( d \) or \( M \). 
This is consistent with the result for covariates standardized via \cref{lemma:units}.
For the unexplained component, the percentage of the parameter space with sign flips is nonzero for even small \( d \) and \( M \), and appears to approach \( 50\% \) as each of \( d \) and \( M \) grow.
This behavior is similar to the result for standardized covariates.
Together, these findings suggest that the conclusions from \cref{proof:explained_fraction,prop:unexplained_fraction} are not simply a result of our choice of units.

\begin{figure}[H]
    \centerline{\includegraphics[width=0.5\linewidth]{Sections/Figures/prob_flip_simulated.pdf}}
    \caption{Percentage of parameter space leading to unexplained component sign flips, without covariate standardization.}
    \label{fig:prob_unexplained_simulated}
\end{figure}

\subsection{Proofs from \cref{sec:reconciling_theory_empirics}}

\subsubsection{Proof of \cref{prop:no_sign_flip}}
\label{proof:no_sign_flip}
Recall that by definition
$$
	U^{(K)} := \mu_H^T \Delta\beta + \Delta\alpha, \quad \quad U^{(H)} := \mu_K^T \Delta\beta + \Delta\alpha.
$$
By assumption, $\sign(\Delta\alpha) = \sign(\mu_H^T\Delta\beta) = \sign(\mu_K^T\Delta\beta)$, and thus $\sign(U^{(H)}) = \sign(U^{(K)})$.

% WTS 1-16-2026: removed this section as it's obvious (leading us to downgrade the proposition to a remark)
%\subsubsection{Proof of \cref{prop:no_explained_sign_flip}}
%\label{proof:no_explained_sign_flip}
%As $\mu_K \neq \mu_H$ by assumption, \cref{lemma:units} applies to imply that we can fix units such that $\mu_K = \mathbf{0}$ and $\mu_H = \mathbf{1}$. Then $\Delta\mu = \mathbf{1}$, and the condition for the explained component to not sign flip, $\sign(\Delta\mu^T\beta_K) \neq \sign(\Delta\mu^T \beta_H)$ reduces to:
%\begin{equation}
%	\sign\left( \sum_{i=1}^d \beta_{K,i} \right) \neq \sign\left( \sum_{i=1}^d \beta_{H,i} \right),
%\end{equation}
%as claimed.

